\Qhead{4.11}{Линейный электрооптический и магнитооптический эффекты}{Линейный электрооптический и магнитооптический эффекты и их применение для управления светом.}
\Subfile{4.11.1}{sections/04_generation/q11/s01}

\recap{Поккельс управляет линейным двулучепреломлением через электрическое поле, Фарадей — невзаимным вращением через магнитное. Для интенсивностной модуляции фазовый эффект дополняют поляризаторами или интерферометром.}
\board{Эллипсоид показателей и $\Delta n$; $V_\pi$; передаточная синусоида; схема изолятора.}
\oral{Почему кристалл Поккельса без анализатора может не менять полную мощность проходящего света?}
\bait{Можно оставить невзаимность Фарадея: при двойном проходе вращение не компенсируется, что позволяет подавлять обратное отражение.}
\sources{Университетские слайды~\cite{optoelectronics}, Поккельс и АО-управление; Warde~\cite{wardeeo}, лабораторная электрооптика; магнитооптическая связь сопоставлена с круговыми модами вопроса~\Qref{2.3}.}
